\chapter{Withdrawn}

\textbf{Purpose:} Record every claim made and withdrawn in the course of assembling this book,
with what killed it, so that a reader can tell which statements elsewhere in these pages survived a
check and which were never tested. \textbf{Scope:} the session of 2026-09-11 that produced this
book, and the two peer sessions that contributed to it.

A withdrawn claim stays on the page with what killed it. An account that records only what survived
tells a reader nothing about how hard the surviving claims were pressed.

Eight claims were withdrawn. Six were this work's own.

\section{The forced case is not the problem}

\textbf{Claimed.} That Fefferman's statements (C) and (D) are not the Millennium problem as the
community means it, and that the prize is reserved for statements (A) and (B), the unforced cases.
The claim was made twice, to two different sessions, and used to argue that the problem remained
open after the September 2026 announcement.

\textbf{What killed it.} Reading Fefferman. He writes that a proof of one of the four statements is
what is asked, and (C) and (D) name a smooth external force inside the statement. The document was
also searched for any passage reserving the prize to (A) or (B); there is none, and the word
\emph{prize} occurs in it once, about the Euler equations.

\textbf{Where it came from.} Journalism, repeated without opening the source. The framing that
working mathematicians exclude forcing from the question they care about is a claim about community
taste, and it was flattened into a claim about the problem statement. The claim had passed through
at least two hands before reaching this work and nobody in that chain had opened the six-page
document that settles it.

\textbf{What it cost to find.} One PDF, six pages, and a search for one word.

\section{The corpus was missing a file}

\textbf{Claimed.} That the corpus was short Fefferman's Navier--Stokes statement, and that the
omission might be an oversight worth correcting.

\textbf{What killed it.} The Clay index of Millennium problems lists five as unsolved, and the
corpus holds exactly those five. It matches the index. The absence was correct and the report of it
as a possible oversight was wrong.

\section{The forced case is out of scope because it is occupied}

\textbf{Claimed.} That the forced case was ``occupied and contested,'' phrased so as to put it
outside what this project should look at.

\textbf{What killed it.} Occupied is a fact and contested is a fact. Out of scope was editorial, and
scope was never this work's to set. A second session accepted the phrase and passed it onward
without asking what it rested on, and recorded afterward that they had applied the
not-a-second-witness rule to the facts in that message while not applying it to the framing.

\section{The Clay page settles it}

\textbf{Claimed}, by the project architect. That the Clay problem page reading \emph{Active} is the
authority on the matter and says the problem is unsolved, and that a primary source would be needed
before accepting otherwise.

\textbf{What killed it.} A webpage lagging an announcement six days old is unremarkable. The claim
leaned on a page's categorization as though it were evidence about the mathematics and not
evidence about a page. Withdrawn by its author on reading the primary source.

\textbf{What survived of it.} The narrower and better half. Whether the formalized statement matches
the problem is the thing to check, and that is where scrutiny belongs. The reading in the
Navier--Stokes chapter exists because this caution was raised.

\section{The uniform null}

\textbf{Claimed}, by the project architect. That distance to the nearest integer is uniform on
$[0,\tfrac12]$ for a real carrying no integrality structure, so agreement to $N$ digits is a bound
of about $2\times10^{-N}$, and fifty digits gives $10^{-50}$.

\textbf{What killed it.} It was derived and not drawn. Withdrawn by its author, who recorded
that seven derived thresholds in one earlier session had every one come out too low and always in
the same direction, because deriving a bar means enumerating the variance sources and the forgotten
ones only add variance.

\textbf{Replaced by} the drawn null in the toolkit chapter, measured through the same pipeline under
a condition where the answer is known to be wrong.

\section{Fifty digits is a bound of ten to the minus fifty}

\textbf{Claimed.} This work relayed the figure above onward as though it were the standard the
measurement would meet, before its author withdrew it.

\textbf{What killed it.} The same withdrawal. Recorded separately because relaying a claim is its
own act: the number reached a third party through this work, and retracting it at the source does
not retract it there.

\section{The integer factors concentrate the quotient}

\textbf{Claimed.} That the integer factors in the Birch and Swinnerton-Dyer quotient, the torsion
order squared, the Tamagawa product and $r!$, might concentrate it near integers and inflate the
chance-coincidence rate, so that the null needed conditioning on the arithmetic structure already
in the inputs.

\textbf{What killed it.} Scaling a structureless real by a known rational does not concentrate it
near integers. The objection was wrong on its mechanism. The conclusion it was offered in support
of, that the null was wrong, happened to be correct for the unrelated and better reason in the
section above.

\section{The architect cited files that do not exist}

\textbf{Claimed}, for the duration of one tool call. That two artifacts cited in the recommendation,
a partition table and a scaling measurement, were absent from the tree.

\textbf{What killed it.} Searching wider. Both paths were repository-relative to BTCminer and had
been resolved against anchor\_sift. Both files exist and every figure quoted from them was
subsequently confirmed in the files.

\textbf{Why it is in this workbook although it was never reported.} Because it was one step from
being reported, and the failure was the same one as the first entry in this chapter: a conclusion
drawn from a check that had not actually been performed against the right object. The near miss is
the useful part.

\section{The pattern}

Of the eight, five were killed by opening a primary source that a summary had flattened. Two were
killed by a peer's framing being accepted without asking what it rested on. One was killed by
resolving a path against the wrong root.

None was killed by a disagreement about mathematics. Every one was killed by reading something that
had been available the whole time.
